\documentclass{fiche}
\metadonnees{5}{Sang-froid}{Bloc 1 — Le récit : présent et passé}

\begin{document}
\entete

\objectifs{%
\textbf{Langue} : bilan des temps du passé ; connecteurs du récit ; thème.\par
\textbf{Texte} : Saki (H.\,H. Munro), \og The Open Window \fg (1914).\par
\textbf{Durée} : 1\,h\,30 — exercices 35 min, texte et questions 25 min, version 30 min.}

%% ===================================================================
\exercice[12 min]{Bilan des temps du passé}

\rappel{\textbf{Rappel.} Prétérit : fait daté, révolu. Present perfect : bilan qui compte
encore maintenant. Pluperfect : fait antérieur à un autre fait passé.}

\consigne{Mettez le verbe entre parenthèses au temps qui convient.}

\begin{anglais}
\begin{enumerate}[label=\arabic*.,itemsep=2.2mm,leftmargin=*]
  \item Framton \emph{(arrive)} \trou{arrived} in the country a month ago.
  \item His sister \emph{(stay)} \trou{had stayed} at the rectory four years before.
  \item ``I \emph{(not meet)} \trou{have not met} a single soul since I came here.''
  \item The three men \emph{(never come)} \trou{never came} back from the marshes.
  \item Nobody \emph{(find)} \trou{has found} their bodies, and nobody ever will.
  \item She \emph{(just tell)} \trou{has just told} him the whole story.
  \item By the time the aunt came in, the niece \emph{(finish)} \trou{had finished} her tale.
  \item ``How long \emph{(you / know)} \trou{have you known} my aunt?'' --- ``Not at all.''
  \item The tide \emph{(be)} \trou{had been} in when they crossed the moor that morning.
  \item Framton \emph{(shiver)} \trou{shivered} slightly and turned towards the girl.
  \item He \emph{(suffer)} \trou{had been suffering} from his nerves for years when the
    doctors finally ordered complete rest.
  \item ``\emph{(you / ever see)} \trou{Have you ever seen} a ghost?'' she asked calmly.
\end{enumerate}
\end{anglais}

\note{Prétérit : 1, 4, 10. Present perfect : 3, 5, 6, 8, 12. Pluperfect : 2, 7, 9, 11.}

\note[Les deux décisions]{D'abord : le fait est-il rattaché à maintenant (present perfect)
ou daté et clos (prétérit) ? Ensuite seulement : est-il antérieur à un autre fait passé
(pluperfect) ? Phrase 4 : \emph{never} n'impose pas le present perfect ; ici les trois hommes
sont morts, l'affaire est close, le prétérit s'impose. Phrase 11 : \emph{had been suffering},
durée antérieure à un moment passé, exactement comme \emph{have been} l'est à maintenant.}

%% ===================================================================
\exercice[8 min]{Connecteurs du récit}

\rappel{\textbf{Rappel.} Les connecteurs de temps ordonnent le récit ; ils ne changent pas
le temps du verbe.}

\consigne{Complétez avec : \emph{then}, \emph{suddenly}, \emph{at last}, \emph{presently},
\emph{meanwhile}, \emph{as soon as}, \emph{by the time}, \emph{while}, \emph{until},
\emph{eventually}.}

\begin{anglais}
\begin{enumerate}[label=\arabic*.,itemsep=2.2mm,leftmargin=*]
  \item \trou{As soon as} the niece had finished, her aunt came into the room.
  \item He waited \trou{until} the window was closed at dusk.
  \item \trou{Suddenly} three figures appeared at the end of the lawn.
  \item \trou{By the time} he reached the gate, the cyclist had swerved into the hedge.
  \item She talked about the shooting; \trou{meanwhile} her eyes never left the window.
  \item \trou{At last} the doctors ordered him complete rest.
  \item He said good morning, and \trou{then} walked out without a word.
  \item \trou{While} he was speaking, the child stared through the window.
  \item \trou{Presently} the spaniel came trotting behind them.
  \item \trou{Eventually} he understood that he had been fooled.
\end{enumerate}
\end{anglais}

\note[Trois faux amis]{\emph{Eventually} ne signifie pas \og éventuellement \fg{} mais
\og finalement, à la longue \fg{} ; \og éventuellement \fg{} se dit \emph{possibly}.
\emph{Presently} ne signifie pas \og présentement \fg{} mais \og bientôt, peu après \fg.
\emph{Actually} ne signifie pas \og actuellement \fg{} mais \og en fait \fg{} ;
\og actuellement \fg{} se dit \emph{currently}.}

\note[Construction]{\emph{Until}, \emph{as soon as}, \emph{by the time} et \emph{while}
introduisent une proposition (sujet + verbe). \emph{Then}, \emph{suddenly}, \emph{at last},
\emph{presently}, \emph{meanwhile}, \emph{eventually} sont des adverbes et se placent en
tête de phrase ou juste avant le verbe.}

%% ===================================================================
\exercice[15 min]{Thème}
\consigne{Traduisez en anglais. Tous les temps du bloc y sont.}

\begin{quote}\itshape
Framton Nuttel n'avait jamais rencontré madame Sappleton. Sa sœur, qui avait habité la
région quatre ans plus tôt, lui avait donné des lettres d'introduction. Depuis un mois, il
rendait visite à des inconnus les uns après les autres, et ses nerfs n'allaient pas mieux.
Autrefois il aimait la campagne ; à présent elle lui faisait peur. Quand la jeune fille lui
demanda s'il connaissait beaucoup de monde dans les environs, il répondit qu'il ne
connaissait pratiquement personne. Il ne savait pas encore que, depuis trois ans, la fenêtre
restait ouverte tous les soirs jusqu'à la nuit tombée.
\end{quote}

\lignes{10}

\corr{\textbf{Proposition.} Framton Nuttel had never met Mrs Sappleton. His sister, who had
stayed in the neighbourhood four years before, had given him letters of introduction. For a
month he had been calling on one stranger after another, and his nerves were no better. He
used to like the country; now it frightened him. When the young lady asked him whether he
knew many people round there, he answered that he hardly knew a soul. He did not yet know
that for three years the window had been left open every evening until it was quite dusk.}

\note[Quatre pièges]{\og quatre ans plus tôt \fg{} : \emph{four years before} et non
\emph{ago} --- \emph{ago} se compte depuis maintenant, \emph{before} depuis un moment passé.
\og Depuis un mois, il rendait visite \fg{} : durée antérieure à un repère passé, donc
\emph{had been calling}. \og Autrefois il aimait \fg{} : \emph{used to like}, pas
\emph{would like}, qui signifierait \og il aimerait \fg. \og demanda s'il \fg{} :
\emph{asked whether} ou \emph{asked if}, jamais \emph{asked that}.}

\note[Deux tournures]{\og rendre visite à quelqu'un \fg{} : \emph{to call on somebody}.
\og il ne connaissait pratiquement personne \fg{} : \emph{he hardly knew a soul} --- une
seule négation en anglais, portée par \emph{hardly}.}

%% ===================================================================
\newpage
\section{Texte}

\noindent\small\textit{Framton Nuttel, qui soigne ses nerfs à la campagne, fait une visite de
politesse chez des gens qu'il ne connaît pas.}\normalsize

\begin{texte}
``My aunt will be down presently, Mr. Nuttel,'' said a very
\gl[self-possessed]{self-possessed}{maître de soi} young lady of fifteen; ``in the meantime
you must try and put up with me.''

Framton Nuttel \gl[to endeavour]{endeavoured}{s'efforcer de} to say the correct something
which should duly flatter the niece of the moment without unduly discounting the aunt that
was to come. Privately he doubted more than ever whether these formal visits on a succession
of total strangers would do much towards helping the nerve cure which he was supposed to be
undergoing.

``I know how it will be,'' his sister had said when he was preparing to migrate to this rural
retreat; ``you will bury yourself down there and not speak to a living soul, and your nerves
will be worse than ever from \gl[to mope]{moping}{broyer du noir}. I shall just give you
letters of introduction to all the people I know there. Some of them, as far as I can
remember, were quite nice.''

Framton wondered whether Mrs. Sappleton, the lady to whom he was presenting one of the
letters of introduction, came into the nice division.

``Do you know many of the people round here?'' asked the niece, when she judged that they had
had sufficient silent communion.

``Hardly a soul,'' said Framton. ``My sister was staying here, at the
\gl[a rectory]{rectory}{le presbytère}, you know, some four years ago, and she gave me
letters of introduction to some of the people here.''

He made the last statement in a tone of distinct regret.

``Then you know practically nothing about my aunt?'' pursued the self-possessed young lady.

``Only her name and address,'' admitted the caller. He was wondering whether Mrs. Sappleton
was in the married or widowed state. An undefinable something about the room seemed to
suggest masculine habitation.

``Her great tragedy happened just three years ago,'' said the child; ``that would be since
your sister's time.''

``Her tragedy?'' asked Framton; somehow in this restful country spot tragedies seemed out of
place.

``You may wonder why we keep that window wide open on an October afternoon,'' said the niece,
indicating a large \gl[a French window]{French window}{une porte-fenêtre} that opened on to a
lawn.

``It is quite warm for the time of the year,'' said Framton; ``but has that window got
anything to do with the tragedy?''

``Out through that window, three years ago to a day, her husband and her two young brothers
went off for their day's shooting. They never came back. In crossing the
\gl[a moor]{moor}{une lande} to their favourite \gl[a snipe]{snipe}{une bécassine}-shooting
ground they were all three engulfed in a treacherous piece of \gl[a bog]{bog}{une
tourbière, un bourbier}. It had been that dreadful wet summer, you know, and places that were
safe in other years \gl[to give way]{gave way}{céder, s'effondrer} suddenly without warning.
Their bodies were never recovered. That was the dreadful part of it.'' Here the child's voice
lost its self-possessed note and became \gl[falteringly]{falteringly}{d'une voix
hésitante} human. ``Poor aunt always thinks that they will come back some day, they and the
little brown \gl[a spaniel]{spaniel}{un épagneul} that was lost with them, and walk in at
that window just as they used to do. That is why the window is kept open every evening till
it is quite dusk. Poor dear aunt, she has often told me how they went out, her husband with
his white waterproof coat over his arm, and Ronnie, her youngest brother, singing `Bertie,
why do you bound?' as he always did to tease her, because she said it got on her nerves. Do
you know, sometimes on still, quiet evenings like this, I almost get a creepy feeling that
they will all walk in through that window ---''

She broke off with a little shudder. It was a relief to Framton when the aunt
\gl[to bustle]{bustled}{entrer en s'affairant} into the room with a whirl of apologies for
being late in making her appearance.

``I hope Vera has been amusing you?'' she said.

``She has been very interesting,'' said Framton.

``I hope you don't mind the open window,'' said Mrs. Sappleton briskly; ``my husband and
brothers will be home directly from shooting, and they always come in this way. They've been
out for snipe in the marshes to-day, so they'll make a fine mess over my poor carpets. So
like you \gl[men-folk]{men-folk}{les hommes, la gent masculine}, isn't it?''

She rattled on cheerfully about the shooting and the scarcity of birds, and the prospects for
duck in the winter. To Framton it was all purely horrible. He made a desperate but only
partially successful effort to turn the talk on to a less \gl[ghastly]{ghastly}{sinistre,
macabre} topic; he was conscious that his hostess was giving him only a fragment of her
attention, and her eyes were constantly straying past him to the open window and the lawn
beyond. It was certainly an unfortunate coincidence that he should have paid his visit on
this tragic anniversary.
\end{texte}
\source{Saki, \og The Open Window \fg, \emph{Beasts and Super-Beasts}, Londres, John Lane,
1914.}

\partiel[12 min]{Questions}
\consigne{\en{Answer in English, in full sentences.}}

\begin{enumerate}[label=\arabic*.,itemsep=2mm,leftmargin=*]
  \item \en{Why is Framton paying these visits, and what does he think of them?}
    \lignes{2}
    \corr{His sister gave him letters of introduction so that he would not bury himself away
    during his nerve cure. He privately doubts that calling on total strangers will do his
    nerves any good.}
  \item \en{What exactly does Framton know about Mrs Sappleton before the niece speaks?}
    \lignes{2}
    \corr{Only her name and address. He does not even know whether she is married or a
    widow.}
  \item \en{The niece asks two questions before telling her story. What is she checking?}
    \lignes{2}
    \corr{That he knows nobody in the neighbourhood and nothing about her aunt, so that
    nobody can contradict her and he cannot check anything.}
  \item \en{Find three details in the niece's story that make it sound true.}
    \lignes{3}
    \corr{The date (\emph{three years ago to a day}), the wet summer and the bog giving way,
    the white waterproof coat, the brown spaniel, the song \emph{Bertie, why do you bound?},
    and her voice breaking at the right moment.}
  \item \en{Mrs Sappleton's first words contradict the niece's story. Why does Framton not
    notice?}
    \lignes{3}
    \corr{He has just been told that the men are dead, so he hears her cheerful talk about
    their return as the talk of a widow who refuses to accept the truth. The straying eyes
    seem to confirm it.}
  \item \en{What is the effect of telling the story from Framton's point of view?}
    \lignes{2}
    \corr{The reader knows only what he knows, and so falls into the same trap. The comedy
    comes afterwards, when both discover they were deceived by the same details.}
\end{enumerate}

%% ===================================================================
\newpage
\section{Version}
\consigne{Traduisez la fin de la nouvelle. Entre-temps, madame Sappleton s'est écriée
\og Les voilà enfin ! \fg{} ; la nièce regarde par la fenêtre, l'horreur dans les yeux.}

\begin{version}
In a chill shock of nameless fear Framton swung round in his seat and looked in the same
direction.

In the deepening twilight three figures were walking across the lawn towards the window; they
all carried guns under their arms, and one of them was additionally burdened with a white
coat hung over his shoulders. A tired brown spaniel kept close at their heels. Noiselessly
they neared the house, and then a hoarse young voice chanted out of the dusk: ``I said,
Bertie, why do you bound?''

Framton grabbed wildly at his stick and hat; the hall-door, the gravel-drive, and the front
gate were dimly-noted stages in his headlong retreat. A cyclist coming along the road had to
run into the hedge to avoid an imminent collision.

``Here we are, my dear,'' said the bearer of the white \gl[a mackintosh]{mackintosh}{un
imperméable}, coming in through the window; ``fairly muddy, but most of it's dry. Who was
that who bolted out as we came up?''

``A most extraordinary man, a Mr. Nuttel,'' said Mrs. Sappleton; ``could only talk about his
illnesses, and dashed off without a word of good-bye or apology when you arrived. One would
think he had seen a ghost.''

``I expect it was the spaniel,'' said the niece calmly; ``he told me he had a horror of dogs.
He was once hunted into a cemetery somewhere on the banks of the Ganges by a pack of
\gl[a pariah dog]{pariah dogs}{un chien errant} and had to spend the night in a newly dug
grave with the creatures snarling and grinning and foaming just above him. Enough to make
anyone lose their nerve.''

Romance at short notice was her speciality.
\end{version}
\source{\emph{Ibid.}}

\lignes{22}

\corr{\textbf{Proposition de traduction.} Saisi d'un frisson glacé, d'une peur sans nom,
Framton pivota sur son siège et regarda dans la même direction.

Dans le crépuscule qui s'épaississait, trois silhouettes traversaient la pelouse vers la
fenêtre ; toutes trois portaient un fusil sous le bras, et l'une d'elles était en outre
chargée d'un manteau blanc jeté sur les épaules. Un épagneul brun, fatigué, les suivait de
près. Sans un bruit, ils approchèrent de la maison, puis une voix jeune et enrouée
psalmodia dans la pénombre : \og Je répète : Bertie, pourquoi bondis-tu ? \fg

Framton empoigna follement sa canne et son chapeau ; la porte du vestibule, l'allée de
gravier et le portail ne furent que des étapes à peine remarquées dans sa fuite éperdue. Un
cycliste qui venait sur la route dut se jeter dans la haie pour éviter la collision.

\og Nous voici, ma chère, dit le porteur de l'imperméable blanc en entrant par la fenêtre ;
un peu crottés, mais c'est presque sec. Qui donc a filé au moment où nous arrivions ? \fg

\og Un homme des plus extraordinaires, un certain monsieur Nuttel, dit madame Sappleton ; il
n'a su parler que de ses maladies et il a détalé sans un mot d'adieu ni d'excuse quand vous
êtes arrivés. On croirait qu'il a vu un fantôme. \fg

\og Ce doit être l'épagneul, dit la nièce avec calme ; il m'a dit qu'il avait les chiens en
horreur. Il a été poursuivi un jour jusque dans un cimetière, quelque part sur les bords du
Gange, par une meute de chiens errants, et il a dû passer la nuit dans une tombe fraîchement
creusée, avec ces bêtes qui grondaient, montraient les dents et écumaient juste au-dessus de
lui. De quoi faire perdre son sang-froid à n'importe qui. \fg

L'affabulation au pied levé était sa spécialité.}

\note[Choix de traduction 1]{\emph{Romance at short notice was her speciality} : \emph{romance}
n'est pas \og la romance \fg{} mais l'invention romanesque ; \emph{at short notice} = \og au
pied levé, sans préavis \fg. Toute la chute tient dans cette phrase : la traduire mot à mot
la détruit.}

\note[Choix de traduction 2]{\emph{one of them was additionally burdened with a white coat
hung over his shoulders} : passif d'état, à rendre par un participe français (\og était en
outre chargée de \fg) et non par \og était été chargé \fg. \emph{Hung over his shoulders}
est un second participe, apposé au manteau.}

\note[Choix de traduction 3]{\emph{Who was that who bolted out as we came up?} :
\emph{to bolt} = \og détaler \fg ; \emph{to come up} = \og s'approcher, arriver \fg. Les
verbes à particule se traduisent par un verbe simple précis, jamais par le verbe de base
plus un adverbe.}

\note[Choix de traduction 4]{\emph{One would think he had seen a ghost} : \emph{one} indéfini
(\og on \fg) et pluperfect après \emph{would think} : \og on croirait qu'il a vu \fg. Le
français remonte d'un cran de moins que l'anglais.}

\note[Choix de traduction 5]{\emph{Enough to make anyone lose their nerve} : \emph{their}
reprend \emph{anyone} --- pluriel de neutralité, courant en anglais, à rendre par un
singulier en français. \emph{To lose one's nerve} = \og perdre son sang-froid \fg, et non
\og perdre ses nerfs \fg.}

%% ===================================================================
\partiel{Lexique à mémoriser}
\begin{itemize}[label=--,itemsep=0.8mm,leftmargin=*]\small
  \item \textbf{self-possessed} — maître de soi ; \textit{calque de} possession de soi
  \item \textbf{to put up with} — supporter, s'accommoder de
  \item \textbf{to endeavour} — s'efforcer de ; \textit{du français} en devoir
  \item \textbf{to mope} — broyer du noir
  \item \textbf{a soul} — une âme ; not a soul \textit{« pas âme qui vive »}
  \item \textbf{a widow} — une veuve ; \textit{masculin} a widower
  \item \textbf{a lawn} — une pelouse
  \item \textbf{dusk} — le crépuscule ; \textit{adjectif} dusky
  \item \textbf{twilight} — la demi-jour, le crépuscule ; twi- \textit{« deux, entre-deux », comme} twice, twin
  \item \textbf{treacherous} — traître, trompeur ; \textit{du français} tricher
  \item \textbf{to give way} — céder, s'effondrer
  \item \textbf{to recover} — retrouver ; \textit{et aussi « se remettre »}
  \item \textbf{a shudder} — un frisson d'horreur
  \item \textbf{creepy} — qui donne la chair de poule ; \textit{de} to creep \textit{« ramper »}
  \item \textbf{to tease} — taquiner
  \item \textbf{briskly} — vivement, d'un ton allègre
  \item \textbf{ghastly} — sinistre, macabre ; \textit{même racine que} ghost
  \item \textbf{to stray} — errer, s'égarer
  \item \textbf{to bolt} — détaler
  \item \textbf{headlong} — éperdu, tête baissée
\end{itemize}

\end{document}
